%% kapa: exam=1 round=27 year=2016 period=2 form=A subject=appraisal_law subject_ko=감정평가관계법규 seq=6 number=6
%% source: Q-Net 감정평가사 기출문제 (cst00309) artlSeq=5169402 file=extracted/1008/5169402/제27회 감정평가사 1차 문제/제27회 감정평가사 2교시문제(A형).pdf
%% conversion: kordoc 4.24.51 -> tools/gen1.py (auto); review=proofread-sonnet; flags=none
\begin{question}{6}
국토의 계획 및 이용에 관한 법령상 도시ㆍ군계획시설 부지의 매수 청구에 관한 설명으로 옳지 않은 것은?
\begin{choices}
  \item 매수청구대상에는 도시․군계획시설의 부지로 되어 있는 토지 중 지목이 잡종지인 토지는 포함되지 않는다.
  \item 부재부동산 소유자의 토지로서 매수대금이 2,000만원을 초과하는 경우 매수의무자는 도시ㆍ군계획시설채권을 발행하여 지급할 수 있다.
  \item 도시ㆍ군계획시설사업의 시행자가 정하여진 경우 매수대상인 토지의 소유자는 시행자에게 그 토지의 매수를 청구할 수 있다.
  \item 도시ㆍ군계획시설채권의 상환기간은 10년 이내로 한다.
  \item 매수의무자가 매수하기로 결정한 토지는 매수 결정을 알린 날부터 2년 이내에 매수하여야 한다.
\end{choices}
\end{question}
\answer{2}
